\scanpage{59}
\begin{tabular}{c@{\qquad}l}
U N I V E R S I T E & MONTPELLIER, le 16 mars 1979\\
DES SCIENCES ET TECHNIQUES &\\
DU LANGUEDOC &\\
$\sim\sim\sim\sim\sim\sim\sim\sim\sim\sim\sim\sim$ &
\end{tabular}

Service de la Scolarité

MR/MG

\uline{Diffusion} : Liste H

N° 1634

\uline{AVIS DE SOUTENANCE D'UNE THESE DE DOCTORAT D'ETAT}

\uline{Nom et Prénom du candidat} : Monsieur LAIGNELET Bernard.

\uline{Titre de la thèse} : "Etude de certains facteurs biochimiques communs à la coloration et à la qualité culinaire des pâtes alimentaires de blé dur".

\uline{Date, heure et lieu de soutenance prévus} : Le 4 avril 1979 à 10 heures, Salle de Cours N° 9 U.S.T.L.

\uline{Nom du Directeur de la thèse} : Monsieur Pierre FEILLET.

\uline{Désignation du laboratoire où la thèse est préparée} : Laboratoire de Technologie des blés durs et du riz -INRA- Montpellier.

\uline{Résumé de la thèse} : Par des analyses statistiques et biochimiques, on montre qu'il n'existe pas de lien entre la coloration des pâtes alimentaires et les propriétés plastiques des pâtes et du gluten de blé dur.
Mais l'évolution de la coloration sous l'effet d'oxydations catalysées par les oxydoréductases présentes dans la semoule limiterait l'oxydation de la pâte et du gluten durant le pétrissage.

\uline{Durée approximative des recherches ayant abouti à la thèse} : 6 ans.

\uline{Noms des rapporteurs} :

M. CROUZET Jean\quad Directeur Scientifique, maître de conférences, Biochimie Alimentai-\\
\hspace*{3cm}re -U.S.T.L.-

M. ESCOUFIER Yves\quad Directeur Scientifique, maître de conférences, Statistique, -U.S.T.L.-

M. FEILLET Pierre\quad Directeur Ingénieur, maître de recherches, Technologie des Céréales -I.N.R A.-

\uline{Composition du jury} :

M. CROUZET J.\quad maître de conférences, Biochimie Alimentaire, U.S.T.L. (Président)

M. DRAPRON R.\quad Ingénieur 2° A, Biochimie Céréales, INRA\\
\hspace*{3cm}(Membre extérieur choisi en raison de sa compétence scientifique)

M. ECOCHARD R.\quad Professeur, Phytotechnie, E.N.S.A.T.

M. ESCOUFIER Y.\quad Maître de conférences, Statistiques, U.S.T.L.\\
\hspace*{3cm}(Membre étranger à la spécialité du candidat).

M. FEILLET P.\quad Maître de recherches, Technologie céréales,I.N.R.A.

Le Secrétaire Général,

MR LE PROFESSEUR GROTHENDIECK ALEX.

\unsure{E. Siau}

MATHEMATIQUES\qquad 5

\uline{E. SIAU}
\mspage{58}
